\title{Controlling Text-to-Image Diffusion by Orthogonal Finetuning}

\begin{abstract}
Recent advancements in large text-to-image diffusion models have enabled highly realistic image synthesis from textual descriptions. However, the question of how to systematically direct or modify these models for specific downstream applications remains a significant challenge. In response, we propose Orthogonal Finetuning~(OFT), a theoretically motivated approach for finetuning text-to-image diffusion models on new tasks. Distinct from previous techniques, OFT demonstrably maintains the hyperspherical energy, a property encapsulating the relationships between neurons on the unit hypersphere. Our analysis indicates that preserving this geometric structure is fundamental for retaining the model's capacity to generate semantically coherent images. To further enhance the robustness of the finetuning process, we introduce Constrained Orthogonal Finetuning (COFT), which incorporates an explicit constraint on the hyperspherical radius. Our work focuses on two representative finetuning scenarios: subject-driven generation—where the aim is to produce images of a particular entity, given limited reference images and a text description—and controllable generation—where additional inputs such as control signals are incorporated into the generative process. Through empirical evaluations, we demonstrate that OFT achieves superior image quality and faster convergence compared to existing methods.
\end{abstract}

\section{Introduction}

State-of-the-art text-to-image diffusion frameworks~\cite{saharia2022photorealistic,ramesh2022hierarchical,rombach2022high} exhibit remarkable success in high-resolution image synthesis steered by textual prompts. Nonetheless, relying solely on textual conditioning often fails to provide the level of nuanced and precise control needed for diverse applications. To mitigate these limitations, this paper investigates two prominent tasks in text-to-image generation:

\begin{itemize}[leftmargin=*,nosep]

    \item \textbf{Subject-driven generation}~\cite{ruiz2023dreambooth}: With only a handful of visual examples of a target subject, the objective is to generate varied images featuring the same entity in novel contexts, guided by textual input.
    \item \textbf{Controllable generation}~\cite{zhang2023adding,mou2023t2i}: When provided with an extra modality, such as edge maps or semantic segmentations, the model is tasked with producing images that not only adhere to these auxiliary signals but also follow textual instructions.
\end{itemize}

Both of these problems center around the challenge of how to fine-tune text-to-image diffusion models so as to retain the generative quality achieved through pretraining. The ideal characteristics of a fine-tuning technique can be summarized as follows: (1) \emph{training efficiency}, which entails using fewer adjustable parameters and minimizing the number of required training epochs, and (2) \emph{generalizability preservation}, meaning the maintenance of the pretrained model's high-fidelity and diverse generative capabilities. In current practice, fine-tuning is generally carried out by either updating the model’s neuron weights with a small learning rate (\emph{e.g.}, \cite{ruiz2023dreambooth}) or by augmenting the model with a lightweight component containing re-parameterized weights (\emph{e.g.}, \cite{hulora2022,zhang2023adding}). Despite their conceptual straightforwardness, these fine-tuning approaches are not sufficient to ensure that the pretrained generative qualities are conserved. Furthermore, there is a notable gap in our theoretical understanding of what constitutes an effective fine-tuning strategy and how to optimally select hyperparameters, such as the training duration. A significant hurdle is the absence of a robust metric for assessing to what extent the original generative capacity is preserved. Current methodologies often presume that minimizing the Euclidean distance between the weights of the fine-tuned and pretrained models corresponds to a higher degree of ability retention. Consequently, these approaches usually adopt very conservative learning rates. Although this intuition can sometimes be correct, we posit that relying solely on the Euclidean proximity to the pretrained model does not sufficiently reflect the degree of semantic conservation. Instead, introducing a measure that captures more structural distinctions between the original and adapted models could substantially enhance both the preservation of pretrained functionality and the overall robustness of fine-tuning.

Building on previous findings that hyperspherical similarity effectively encapsulates semantic relationships~\cite{liu2017deep,liu2018decoupled,chen2020angular}, we utilize hyperspherical energy~\cite{liu2018learning} as a metric to capture the pairwise interactions among neurons within a layer. Hyperspherical energy represents the aggregate hyperspherical similarities (\emph{e.g.}, cosine similarities) between all pairs of neurons in a given layer, thus reflecting how uniformly neurons are distributed on the unit hypersphere~\cite{liu2021learning}. We propose that, for optimal finetuning, the difference in hyperspherical energy between the pretrained and finetuned models should be minimal. A straightforward strategy would be to introduce a regularization term to keep the hyperspherical energy constant throughout finetuning, yet this does not inherently ensure effective minimization of energy discrepancies. Consequently, we exploit a crucial property of hyperspherical energy invariance: if all neurons undergo an identical orthogonal transformation, their pairwise hyperspherical similarity remains unchanged. Guided by this invariance, we introduce Orthogonal Finetuning~(OFT), a methodology designed to adapt large text-to-image diffusion models for downstream tasks without altering the original hyperspherical energy. The principal mechanism is to learn a shared orthogonal transformation across neurons in a layer, preserving their relative angles. OFT can alternatively be interpreted as reorienting the coordinate system for neurons within a layer. Furthermore, acknowledging that reduced Euclidean distance between the finetuned and pretrained models correlates with superior retention of pretraining capabilities, we propose a refined version, Constrained Orthogonal Finetuning~(COFT), which limits the finetuned model to remain within a hypersphere of constant radius centered at the pretrained neurons.

The rationale behind the effectiveness of orthogonal transformations in neuron finetuning is partly informed by principles from the 2D Fourier transform. Using this mathematical tool, an image is separated into its magnitude and phase spectra, where the phase spectrum—encoding the angular relationships between the input and the basis—carries the predominant semantic information. Notably, reconstructing an image with its original phase spectrum combined with an arbitrary magnitude spectrum often retains the core semantics. This observation implies that adjusting neuron directions is central to achieving meaningful semantic alterations in generated images, aligning with the objectives of both subject-driven and controllable image generation. Nevertheless, allowing unrestrained freedom in changing neuron directions can compromise the generative capabilities established during pretraining. To restrict this freedom, we advocate for the conservation of the inter-neuron angles, underpinned by the assumption that these angles are integral to how neural networks encode information. Building upon this, we introduce the idea of learning a layer-wide orthogonal transformation for neurons, thereby maintaining constant hyperspherical energy throughout each layer.

Our approach is further informed by prior work in orthogonal over-parameterized training~\cite{liu2021orthogonal}, where classification networks are trained from their random initializations via orthogonal transformations. The foundation here is that a randomly initialized network possesses low expected hyperspherical energy, and the strategy in \cite{liu2021orthogonal} is to keep this energy low over the course of training—an approach linked to superior generalization in classification~\cite{liu2018learning,lin2020regularizing}. That work demonstrated the sufficiency of orthogonal transformations for developing highly generalizable classifiers. In contrast, our focus shifts to the domain of finetuning text-to-image diffusion models, targeting improvements in controllability as well as in downstream generative fidelity. The Orthogonal Finetuning Transformation (OFT) we propose is distinct from \cite{liu2021orthogonal} in two key ways. First, where \cite{liu2021orthogonal} aims to reduce hyperspherical energy, OFT is designed to retain the hyperspherical energy intrinsic to the pretrained model, safeguarding its internal structure against potential degradation during finetuning—a necessary consideration since lowering hyperspherical energy during diffusion model finetuning risks erasing existing semantic structure. Second, OFT explicitly minimizes divergence from the pretrained model, introducing a constrained optimization aspect not present in \cite{liu2021orthogonal}. Finetuning thus requires balancing adaptability with the preservation of foundational knowledge; we contend that the OFT framework successfully navigates this balance. Our main contributions are summarized as follows:

\begin{itemize}[leftmargin=*,nosep]

    \item We introduce a new finetuning technique, Orthogonal Finetuning (OFT), designed to enhance the controllability of text-to-image diffusion models. To bolster stability, we further propose a constrained version of OFT that restricts the angular deviation from the pretrained model parameters.
    \item In contrast to other existing finetuning strategies, OFT updates model parameters in a way that rigorously preserves hyperspherical energy, a quantity we empirically demonstrate to correlate with semantic integrity in the pretrained generative model.
    \item The OFT methodology is evaluated on two applications: subject-driven and controllable generation. Extensive experiments reveal that OFT not only yields notable gains in generation quality, finetuning stability, and convergence speed compared to prior methods, but also offers superior data efficiency, achieving robust convergence with fewer training samples and epochs.
    \item For the task of controllable generation, we devise a novel control consistency metric, designed to assess model controllability by extracting the control signal from the generated image and quantitatively comparing it to the original input control signal. This metric underscores the enhanced controllability delivered by OFT.
\end{itemize}

\section{Related Work}

\textbf{Text-to-image diffusion models}. The field of text-to-image synthesis has witnessed remarkable advances~\cite{saharia2022photorealistic,ramesh2022hierarchical,rombach2022high,gu2022vector,nichol2022glide} due to breakthroughs in diffusion-based generative models~\cite{ho2020denoising,song2021score,song2021denoising,dhariwal2021diffusion} and joint vision-language representation learning~\cite{su2020vlbert,lu2019vilbert,radford2021learning,wang2021simvlm,li2022blip,li2023blip,alayrac2022flamingo,schuhmann2022laion}. GLIDE~\cite{nichol2022glide} and Imagen~\cite{saharia2022photorealistic} train diffusion processes directly in pixel space; GLIDE employs end-to-end text encoder training alongside a diffusion prior based on paired datasets, whereas Imagen leverages a fixed, pretrained text encoder. In comparison, latent-space approaches such as Stable Diffusion~\cite{rombach2022high} and DALL-E2~\cite{ramesh2022hierarchical} utilize encoded visual representations; Stable Diffusion employs a VQ-GAN~\cite{esser2021taming} visual codebook as its latent domain, while DALL-E2 incorporates CLIP~\cite{radford2021learning} to unify image and text embeddings. Beyond diffusion models, generative adversarial networks~\cite{reed2016generative,zhang2017stackgan,xu2018attngan,li2019controllable} and autoregressive models~\cite{ramesh2021zero,ding2021cogview,wu2022nuwa,yuscaling} have also been explored for text-conditional image generation. OFT is not tied to a specific architecture, and is broadly applicable to any text-to-image diffusion model.

\textbf{Subject-driven generation}. Existing techniques for mitigating subject modification include using user-supplied masks as additional guidance~\cite{avrahami2022blended,nichol2022glide}. Alternatively, inversion-based approaches~\cite{choi2021ilvr,dhariwal2021diffusion,ramesh2022hierarchical,gal2022image} allow modification of background or context without altering the primary subject. The work by~\cite{hertz2022prompt} enables both localized and global editing in the absence of explicit masks. However, such strategies often fall short in preserving detailed identity characteristics. Pivotal Tuning~\cite{roich2022pivotal} addresses this limitation by finetuning generative models around an initial inverted code with extra identity preservation regularization. Analogously,~\cite{nitzan2022mystyle} devises personalized facial priors from curated datasets of an individual's facial images. The approach of~\cite{casanova2021instance} generates various depictions of an instance, but may compromise on the instance-specific details. DreamBooth~\cite{ruiz2023dreambooth} leverages a customized token and a handful of subject photographs to finetune text-to-image diffusion models, employing a reconstruction loss along with a class-based prior preservation objective. Building upon this, OFT follows the DreamBooth paradigm but replaces standard small-step finetuning with finetuning via orthogonal transformations.

\textbf{Controllable generation}. Image-to-image translation is typically considered a form of controllable generation. Earlier approaches often leverage conditional generative adversarial networks~\cite{isola2017image,zhu2017toward,wang2018high,park2019semantic,choi2018stargan} to tackle this problem. Diffusion models have recently been utilized for this task as well~\cite{saharia2022palette,voynov2022sketch,wang2022pretraining}. Among the latest advancements, ControlNet~\cite{zhang2023adding} introduces a strategy to guide pretrained diffusion models by finetuning them with supplementary control cues, achieving notable performance in controllability. Similarly, T2I-Adapter~\cite{mou2023t2i} concurrently enhances control in generated imagery by finetuning diffusion models. Building on the task frameworks introduced in~\cite{zhang2023adding,mou2023t2i}, our method applies OFT for finetuning diffusion models, consistently enhancing control with fewer samples and reduced finetuning overhead. Notably, OFT incurs no extra computational cost during inference.

\textbf{Model finetuning}. There is an increasing trend towards finetuning large pretrained models for downstream applications~\cite{devlin2018bert,brown2020language,he2022masked}. Adaptation strategies, such as those explored in~\cite{hulora2022,houlsby2019parameter,pfeiffer2020adapterfusion}, are extensively investigated in natural language processing. Of particular relevance to OFT is LoRA~\cite{hulora2022}, which assumes the parameter updates during finetuning are low-rank and additive. By contrast, OFT introduces a layer-wise shared orthogonal transformation that modifies neuron weights multiplicatively and rigorously maintains the angular relationships among neurons within a layer, resulting in improved training stability.

\section{Orthogonal Finetuning}

\subsection{Why Does Orthogonal Transformation Make Sense?}

To motivate orthogonal transformation in the context of finetuning text-to-image diffusion models, we break the reasoning into two parts: (1) the importance of adjusting neuron angles (i.e., their directional component) during finetuning, and (2) the rationale for employing orthogonal transformations for this angular adjustment.

Addressing the first question, we are motivated by empirical findings from \cite{liu2018decoupled,chen2020angular}, which indicate that differences in the angular components of features serve as effective indicators of semantic disparity. SphereNet~\cite{liu2017deep} provides evidence that constraining neuron outputs to lie on a unit hypersphere not only preserves model capacity but also improves generalization, suggesting that the directional attributes of neurons are sufficient to encapsulate the core data semantics. To illustrate the significance of neuron orientations, we devise a simplified experiment as presented in Figure~\ref{fig:toy_ae}, where a conventional convolutional autoencoder is trained on images of flowers. During training, standard inner products are employed to form the feature map ($z$ represents the convolution output for kernel $\bm{w}$ with input patch $\bm{x}$). At test time, we assess three distinct strategies for constructing the feature map: (a) the inner product, mirroring the training phase, (b) exclusively utilizing magnitude information, and (c) relying solely on angular information. According to Figure~\ref{fig:toy_ae}, leveraging the angular aspects of neurons nearly reconstructs the original images, whereas magnitude information fails to retain useful content. It is notable that cosine-based activation (method c) is only applied during inference and is not present during training, which uses purely inner product-based computation. These findings point to neuron directions as the principal repository of semantic details in the input images, implying that altering the orientation of neurons could be more impactful for modifying semantic content.

Regarding the second question, the most straightforward method for refining neuron directions is to collectively apply a rotation or reflection to all neurons within a layer, thus inherently involving an orthogonal transformation. Although adopting more adaptable angular transformations that assign distinct rotations to individual neurons is conceivable, orthogonal transformations emerge as a favorable compromise between adaptability and constraint. Additionally, the work of \cite{liu2021orthogonal} demonstrates that orthogonal transformations suffice for effective neural network learning. To substantiate this claim, we conduct a regularization experiment using the ControlNet~\cite{zhang2023adding} framework for controllable generation. Figure~\ref{fig:ortho} summarizes the outcomes: the conventional OFT method consistently achieves robust and precise control after 20 epochs, while omitting the orthogonality constraint in OFT leads to failures in generating plausible images and yields no controllability. This experiment corroborates the necessity of orthogonal constraints within OFT.

\subsection{General Framework}

Traditional finetuning procedures often employ gradient descent with a reduced learning rate to adjust model parameters (or particular model layers). The use of a small learning rate acts as a soft constraint, limiting the updates so that the adjusted model remains close to the pretrained counterpart; standard finetuning can thus be interpreted as seeking to minimize $\thickmuskip=2mu \medmuskip=2mu \| \bm{M} - \bm{M}^0\|$, where $\bm{M}$ and $\bm{M}^0$ denote the finetuned and original pretrained weights, respectively. While intuitively reasonable, this implicit restriction may still allow excessive flexibility when adapting large-scale models. To enforce stronger constraints, LoRA augments the update process with a low-rank limitation, specifically imposing $\thickmuskip=2mu \medmuskip=2mu \text{rank}(\bm{M} - \bm{M}^0)=r'$, with $r'$ chosen as a small integer. In contrast, OFT adopts a different approach by constraining the pairwise similarities among neurons, formalized as $\thickmuskip=2mu \medmuskip=2mu \| \text{HE}(\bm{M})-\text{HE}(\bm{M}^0) \|=0$, where the hyperspherical energy, $\text{HE}(\cdot)$, is computed for the weight matrix. To clarify this, consider a dense (fully connected) layer $\thickmuskip=2mu \medmuskip=2mu \bm{W}=\{\bm{w}_1,\cdots,\bm{w}_n\}\in\mathbb{R}^{d\times n}$, where each neuron $\bm{w}_i$ lies in $\mathbb{R}^{d}$ ($\bm{W}^0$ corresponds to the pretrained weights). Given input $\thickmuskip=2mu \medmuskip=2mu \bm{x}\in\mathbb{R}^d$, the output is computed as $\thickmuskip=2mu \medmuskip=2mu \bm{z}=\bm{W}^\top\bm{x}$ with $\bm{z}\in\mathbb{R}^n$. The essence of OFT is to minimize the gap in hyperspherical energy between the finetuned and the original model, expressed as:

\begin{equation}\label{eq:oft_principle}
\footnotesize
\min_{\bm{W}} \left\| \text{HE}(\bm{W})-\text{HE}(\bm{W}^0)\right\|~~\Leftrightarrow~~\min_{\bm{W}} \bigg{\|} \sum_{i\neq j} \|\hat{\bm{w}}_i-\hat{\bm{w}}_j\|^{-1}-\sum_{i\neq j} \|\hat{\bm{w}}^0_i-\hat{\bm{w}}^0_j\|^{-1}\bigg{\|}
\end{equation}

where $\thickmuskip=2mu \medmuskip=2mu \hat{\bm{w}}_i:=\bm{w}_i/\|\bm{w}_i\|$ denotes the normalized form of the $i$-th neuron, and hyperspherical energy for a weight matrix $\bm{W}$ is $\thickmuskip=2mu \medmuskip=2mu \text{HE}(\bm{W}):= \sum_{i\neq j} \|\hat{\bm{w}}_i-\hat{\bm{w}}_j\|^{-1}$. It is straightforward to see that the minimum value attainable for Eq.~\eqref{eq:oft_principle} is zero. This minimum is realized provided $\bm{W}$ can be obtained from $\bm{W}^0$ by a rotation or reflection, i.e., when $\thickmuskip=2mu \medmuskip=2mu \bm{W}=\bm{R}\bm{W}^0$ for an orthogonal matrix $\thickmuskip=2mu \medmuskip=2mu \bm{R}\in\mathbb{R}^{d\times d}$ (with determinant $1$ for rotation and $-1$ for reflection). This approach underlies the OFT method: rather than updating arbitrary parameters, OFT restricts finetuning to learning a shared orthogonal matrix for each layer to transform its neurons. More precisely, given a pretrained fully

\begin{equation}\label{eq:oft_general}
    \footnotesize
    \bm{z} = \bm{W}^\top \bm{x} = (\bm{R} \cdot \bm{W}^0)^\top \bm{x},~~~\text{s.t.}~\bm{R}^\top \bm{R} = \bm{R}\bm{R}^\top = \bm{I}
\end{equation}

Here, $\bm{W}$ corresponds to the weight matrix adjusted via OFT-finetuning, and $\bm{I}$ represents the identity matrix. Figure~\ref{fig:block} offers a depiction of OFT. In line with the zero initialization strategy of LoRA, it is necessary to guarantee that OFT commences finetuning directly from $\bm{W}^0$—the original pretrained parameters. To this end, $\bm{R}$, the orthogonal transformation matrix, is initialized as the identity so that the model's initial weights remain unchanged. Maintaining the orthogonality of $\bm{R}$ can be achieved by applying differentiable orthogonalization techniques as suggested by \cite{liu2021orthogonal,lezcano2019cheap}. We will elaborate on practical methods for ensuring this orthogonality in an efficient manner.

\subsection{Efficient Orthogonal Parameterization}\label{sect:eff_ortho}

Common approaches to orthogonalization, such as the differentiable Gram-Schmidt procedure, are generally computationally intensive in practice~\cite{liu2021orthogonal}. To improve efficiency, we utilize Cayley parameterization for generating the orthogonal matrix. More specifically, the orthogonal matrix is constructed as $\thickmuskip=2mu \medmuskip=2mu \bm{R} = (\bm{I} + \bm{Q})(I - \bm{Q})^{-1}$, where $\bm{Q}$ is skew-symmetric, satisfying $\thickmuskip=2mu \medmuskip=2mu \bm{Q} = -\bm{Q}^\top$. This approach provides computational benefits but does impose a restriction: Cayley parameterization yields only orthogonal matrices with determinant $1$—a subset known as the special orthogonal group. Nonetheless, our empirical results show that this restriction does not impair practical model performance. Even with Cayley transforms, however, direct parameterization of $\bm{R}$ becomes inefficient for large $d$. To overcome this, we advocate representing $\bm{R}$ using a block-diagonal structure composed of $r$ blocks, which can be written as:

\begin{equation}
    \footnotesize
    \bm{R} = \text{diag}(\bm{R}_1, \bm{R}_2, \cdots, \bm{R}_r) = \begin{bmatrix}
    \bm{R}_1 \in \text{O}(\frac{d}{r}) &  &\\
    &\ddots & \\
    & & \bm{R}_r \in \text{O}(\frac{d}{r})
    \end{bmatrix} \in \text{O}(d)
\end{equation}

where $\text{O}(d)$ represents the orthogonal group in $d$ dimensions, and $\thickmuskip=2mu \medmuskip=2mu \bm{R} \in \mathbb{R}^{d\times d}$ as well as $\thickmuskip=2mu \medmuskip=2mu \bm{R}_i \in \mathbb{R}^{d/r \times d/r},\forall i$ denote orthogonal matrices. If $\thickmuskip=2mu \medmuskip=2mu r=1$, the block-diagonal orthogonal matrix degenerates into a conventional full orthogonal matrix without constraints. A standard orthogonal matrix of shape $\thickmuskip=2mu \medmuskip=2mu d \times d$ possesses $\thickmuskip=2mu \medmuskip=2mu d(d-1)/2$ free parameters, implying parameter complexity $\mathcal{O}(d^2)$. For an orthogonal matrix structured as an $r$-block diagonal, the number of independent parameters is reduced to $\thickmuskip=2mu \medmuskip=2mu d(d/r-1)/2$, with complexity scaling as $\thickmuskip=2mu \medmuskip=2mu \mathcal{O}(d^2/r)$. Parameter sharing can be further leveraged by setting $\thickmuskip=2mu \medmuskip=2mu \bm{R}_i = \bm{R}_j, \forall i\neq j$, further decreasing the parameter count to $\mathcal{O}(d^2/r^2)$. Importantly, all of these methods maintain the orthogonality of the resulting matrix $\bm{R}$, ensuring that hyperspherical energy is still fully preserved.

We now compare the parameter efficiency of OFT with LoRA. For LoRA, when the low-rank parameter is $r'$, the number of trainable parameters is $\thickmuskip=2mu \medmuskip=2mu r'(d+n)$. If both $r$ and $r'$ scale with the neuron dimension $d$ (for instance, let $\thickmuskip=2mu \medmuskip=2mu r=r'=\alpha d$ for a constant $\thickmuskip=2mu \medmuskip=2mu 0<\alpha\leq 1$), then the trainable parameter count for LoRA is $\mathcal{O}(d^2+dn)$, whereas OFT requires only $\mathcal{O}(d)$. For illustration, take a weight matrix of size $\thickmuskip=2mu \medmuskip=2mu 128\times 128$; with $\thickmuskip=2mu \medmuskip=2mu r'=8$, LoRA yields $2,048$ parameters, while for $\thickmuskip=2mu \medmuskip=2mu r=8$ and no shared blocks, OFT uses $960$ trainable parameters.

\subsection{Constrained Orthogonal Finetuning}

OFT’s adaptability can be curtailed by explicitly bounding the range within which the finetuned model deviates from the pretrained weights. Constrained Orthogonal Finetuning (COFT) achieves this by imposing a restriction such that the fine-tuned model remains close to the original in parameter space. The model’s forward pass is defined by:

\begin{equation}\label{eq:coft_general}
    \footnotesize
    \bm{z}=\bm{W}^\top\bm{x}=(\bm{R}\cdot\bm{W}^0)^\top\bm{x},~~~\text{s.t.}~\bm{R}^\top\bm{R}=\bm{R}\bm{R}^\top=\bm{I},~\norm{\bm{R}-\bm{I}}\leq \epsilon
\end{equation}

where two constraints are imposed: orthogonality and a maximum allowable deviation $\epsilon$ from the identity matrix. While Section~\ref{sect:eff_ortho} introduces the Cayley parameterization as a means to guarantee orthogonality, integrating the $\epsilon$-neighborhood constraint into the Cayley representation is challenging. To probe the behavior of the Cayley transform, we approximate $\thickmuskip=2mu \medmuskip=2mu \bm{R}=(\bm{I}+\bm{Q})(I-\bm{Q})^{-1}$ using the Neumann series,

series converges in the operator norm). Consequently, we are able to reformulate the constraint $\thickmuskip=2mu \medmuskip=2mu\|\bm{R}-\bm{I}\|\leq\epsilon$ within the Cayley transform framework, leading to an equivalent formulation $\thickmuskip=2mu \medmuskip=2mu \|\bm{Q}-\bm{0}\|\leq\epsilon'$ where $\bm{0}$ denotes the zero matrix, and $\epsilon'$ is an alternative error tolerance distinct from $\epsilon$. This reparameterized constraint on $\bm{Q}$ lends itself well to enforcement via projected gradient descent. For initialization that yields an identity mapping for the orthogonal matrix $\bm{R}$, we start $\bm{Q}$ at zero. The COFT method may thus be regarded as imposing two clear constraints: minimal change in hyperspherical energy and a controlled divergence from the pretrained parameters. Whereas the latter is often an implicit regularization in many existing finetuning approaches, COFT brings this aspect into the foreground by making it explicit. While OFT achieves strong performance, we find that COFT further enhances finetuning stability because of this overt constraint on deviation. Figure~\ref{fig:eps_coft} illustrates the sensitivity of COFT to the hyperparameter $\epsilon$. The results reveal that $\epsilon$ governs how much the finetuning process may depart from the original model. For large $\epsilon$, the behavior of COFT approaches that of OFT. As $\epsilon$ decreases, the outputs from COFT increasingly match those from the original pretrained text-to-image diffusion model.

\subsection{Re-scaled Orthogonal Finetuning}

We introduce an augmentation to standard OFT by learning a scaling factor for each neuron, inspired by the observation that neuron rescaling does not affect hyperspherical energy (since magnitudes are normalized out). Concretely, the forward computation adopts the form $\thickmuskip=2mu \medmuskip=2mu\bm{z}=(\bm{R}\bm{W}^0\bm{D})^\top\bm{x}$\footnote{\textbf{Errata}: The NeurIPS camera-ready paper mistakenly reports the re-scaled OFT forward pass as $\thickmuskip=2mu \medmuskip=2mu\bm{z}=(\bm{D}\bm{R}\bm{W}^0)^\top\bm{x}$, though the original code and results in Appendix~\ref{app:rescaled} are correct.}, where $\thickmuskip=2mu \medmuskip=2mu \bm{D}=\text{diag}(s_1,\cdots,s_n)\in\mathbb{R}^{n\times n}$ is a trainable diagonal matrix, with all $s_1,\cdots,s_n>0$. This adjustment makes both $\bm{D}$ and the orthogonal matrix $\bm{R}$ parameters to be optimized, as opposed to only $\bm{R}$ in the original OFT (cf. Eq.~\eqref{eq:oft_general}). This re-scaled variant brings additional flexibility with a negligible increase in parameter count. In our experimental evaluations, we employ the original OFT to highlight the impact of the orthogonal transformation alone, but note that the re-scaled version usually brings further improvements (see Appendix~\ref{app:rescaled}).

\section{Intriguing Insights and Discussions}

\textbf{OFT applies across model architectures}. In principle, OFT is a general approach and can be leveraged with a wide range of neural network designs. In Transformer models, LoRA modifications are most commonly employed on the attention weights~\cite{hulora2022}. To enable a direct comparison with LoRA, our experiments confine the use of OFT to finetuning attention weights. OFT is also advantageous for adapting convolutional layers; specifically, the block-diagonal form of $\bm{R}$ acquires meaningful interpretations in the convolutional setting—an attribute not shared by LoRA. When the number of blocks matches the number of input channels, each block exclusively transforms a particular neuron channel, analogous to tuning depth-wise convolution filters~\cite{chollet2017xception}. If all blocks in $\bm{R}$ are tied, the method reduces to a shared orthogonal transformation for all channels. We offer an initial exploration of OFT-based finetuning for convolutional layers in Appendix~\ref{app:convolution}.

\textbf{Connection to LoRA}. LoRA incorporates a low-rank matrix addition to prevent substantial shifts in the information captured by the pretrained weight matrix. By comparison, OFT enforces that the transformation applied to the pretrained weights remains orthogonal and full-rank, thereby preserving the integrity of the pretrained model's information. The forward computation in OFT can be reformulated as $\thickmuskip=2mu \medmuskip=2mu \bm{z}=(\bm{R}\bm{W}^0)^\top\bm{x}=(\bm{W}^0+(\bm{R}-\bm{I})\bm{W}^0)^\top\bm{x}$, with $\thickmuskip=2mu \medmuskip=2mu (\bm{R}-\bm{I})\bm{W}^0$ functioning similarly to the low-rank adjustment in LoRA. Since $\bm{W}^0$ typically has full rank, OFT similarly introduces low-rank updates when $\thickmuskip=2mu \medmuskip=2mu \bm{R}-\bm{I}$ is of low rank. Like LoRA’s rank hyperparameter $r'$, OFT employs a diagonal block size $r$ to constrain the number of parameters subject to training. Notably, LoRA and OFT embody fundamentally different strategies for achieving parameter efficiency: LoRA leverages low-rank representations, whereas OFT relies on exploiting structural sparsity, such as block-diagonal orthogonality.

\textbf{Why OFT converges faster}. The reasoning is twofold: First, as illustrated in Figure~\ref{fig:toy_ae}, adjusting the directionality of neurons is most effective for semantic modification, aligning precisely with OFT’s mechanism. Second, OFT can be interpreted as tuning neuron parameters along a smooth hypersphere, which is associated with a more favorable optimization geometry. Empirical results supporting this observation can also be found in \cite{liu2021orthogonal}.

\textbf{Why not minimize hyperspherical energy}. Unlike \cite{liu2021orthogonal}, we deliberately avoid the goal of minimizing hyperspherical energy. In classification contexts, neuron representations are most useful when non-redundant. Achieving minimum hyperspherical energy results in a uniform distribution of neurons across the hypersphere, but this objective is ill-suited for finetuning because it risks erasing crucial information from the pretrained weights.

\textbf{Balancing flexibility and regularity in finetuning}. Our findings highlight a fundamental trade-off between model adaptability and constraint. Although standard finetuning offers maximal flexibility, it often suffers from instability and is susceptible to model collapse. In contrast, OFT, despite its straightforwardness, effectively balances adaptability and regularization by maintaining pairwise neuron angular relationships. Furthermore, the block-diagonal approach acts as an enhanced regularization on the orthogonal transformation matrix.

\textbf{No increase in inference complexity}. In distinction to ControlNet, our OFT methodology does not impose any extra computational load during inference on the adapted model. At inference, the learned orthogonal matrix $\bm{R}$ is directly multiplied with the original weight matrix $\bm{W}^0$, resulting in an updated weight matrix $\thickmuskip=2mu \medmuskip=2mu \bm{W}=\bm{R}\bm{W}^0$, thus maintaining inference speed identical to that of the original pretrained model.

\section{Experiments and Results}

\textbf{Overall experimental configuration}. For our experiments, we employ Stable Diffusion v1.5~\cite{rombach2022high} as the baseline text-to-image pretrained model. To ensure unbiased evaluation, image samples from all compared methods are randomly selected. For subject-driven tasks, we follow the experimental protocols of DreamBooth~\cite{ruiz2023dreambooth}, and for controllable generation, the procedures are based on ControlNet~\cite{zhang2023adding} and T2I-Adapter~\cite{mou2023t2i}. In order to directly compare against LoRA, OFT or COFT is always applied to the same layer as LoRA. Supplementary experiments and implementation details are reported in Appendix~\ref{app:settings}.

\subsection{Subject-driven Generation}

\textbf{Experimental setup}. Our baseline methods are DreamBooth~\cite{ruiz2023dreambooth} and LoRA~\cite{hulora2022}. All methods employ the loss function specified in DreamBooth. We adhere to the hyperparameter choices suggested by the respective original papers for both DreamBooth and LoRA to ensure optimal performance. Additional experimental outcomes can be found in Appendix~\ref{app:settings},\ref{app:coft_oft},\ref{app:more_qualitative},\ref{app:failure}.

\textbf{Finetuning stability and convergence}. To begin, we assess the stability of finetuning procedures and the convergence rates for DreamBooth, LoRA, OFT, and COFT, as illustrated in Figure~\ref{fig:teaser} and Figure~\ref{fig:db_convergence}. Our findings indicate that COFT and OFT facilitate highly stable finetuning of the diffusion model. At the 400-iteration mark, DreamBooth and both OFT-based methods exhibit strong subject control, whereas LoRA is unsuccessful at maintaining subject identity. With 2000 iterations, DreamBooth encounters issues with image collapse, and LoRA continues to misrepresent the intended attribute (generating yellow fur instead of a yellow shirt). By contrast, OFT and COFT continue to exercise reliable and consistent control over image synthesis throughout training. These outcomes demonstrate that the OFT framework offers swift yet robust convergence for subject-driven image generation. Additionally, we observe that both OFT and COFT exhibit low sensitivity to the specific number of finetuning iterations, significantly simplifying hyperparameter selection across varying subjects. Consequently, one may assign a large iteration count for OFT and COFT without the need for meticulous adjustment. In the case of COFT, provided an appropriate $\epsilon$ is chosen, both the learning rate and total iterations require minimal tuning effort.

\textbf{Quantitative comparison}. In line with the evaluation protocol of \cite{ruiz2023dreambooth}, we carry out a quantitative analysis targeting subject fidelity (using DINO~\cite{caron2021emerging} and CLIP-I~\cite{radford2021learning}), text prompt adherence (CLIP-T~\cite{radford2021learning}), and sample diversity (LPIPS~\cite{zhang2018unreasonable}). Specifically, CLIP-I involves measuring the average pairwise cosine similarity between CLIP embeddings derived from generated and reference images; DINO is computed analogously, utilizing ViT S/16 DINO embeddings in place of CLIP representations. CLIP-T captures the mean cosine similarity between the embeddings of the provided text prompt and those of the generated images. We further measure the average LPIPS cosine similarity across generated images that correspond to the same subject and prompt. According to Table~\ref{table:dreambooth_final}, both COFT and OFT show a marked improvement over DreamBooth and LoRA on the DINO and CLIP-I metrics, and deliver marginally superior or on-par results for text prompt consistency and diversity. Each approach is evaluated by repeatedly finetuning an identical pretrained model 30 times with distinct random seeds to minimize stochastic effects. The evidence thus points to OFT as delivering not only enhanced convergence speed and reliability, but also consistently superior end-task performance.

\textbf{Qualitative comparison}. To provide a clearer perspective on the advantages brought by OFT, we present a selection of randomly chosen subject-driven generation outcomes in Figure~\ref{fig:dreambooth_final}. In the interest of fairness, each sample is derived from an identical finetuned model per method, ensuring that no text prompt receives individualized optimization tailored to its final result. For each technique, the model with the optimal validation CLIP score is utilized. As observed in Figure~\ref{fig:dreambooth_final}, both OFT and COFT consistently maintain strong semantic fidelity to the subject, whereas LoRA frequently struggles to retain subject identity (\emph{e.g.}, in the bowl case, LoRA entirely misses the subject identity). Furthermore, both OFT and COFT demonstrate substantially improved precision in responding to textual prompts, while DreamBooth, although generally effective at preserving subject characteristics, regularly does not succeed in accurately generating content consistent with the text prompt (\emph{e.g.}, the first row for the bowl scenario). These qualitative results suggest that OFT achieves a superior balance of subject preservation and prompt-based controllability. Additionally, OFT’s performance shows robustness to the number of training iterations, maintaining effectiveness even with extensive iterations, a property not shared by DreamBooth or LoRA. Additional qualitative illustrations can be found in Appendix~\ref{app:more_qualitative}. Human evaluation results, described in Appendix~\ref{app:human}, further corroborate OFT’s superiority.

\subsection{Controllable Generation}

\textbf{Settings}. As comparative baselines, we incorporate ControlNet~\cite{zhang2023adding}, T2I-Adapter~\cite{mou2023t2i}, and LoRA~\cite{hulora2022}. Our experiments span three complex controllable generation tasks, specifically: Canny edge to image~(C2I) using the COCO dataset~\cite{lin2014microsoft}, segmentation map to image~(S2I) on the ADE20K dataset~\cite{zhou2017scene}, and landmark to face~(L2F) based on the CelebA-HQ dataset~\cite{karras2018progressive,xia2021tedigan}. Each method is applied to finetune Stable Diffusion~(SD) v1.5 across these datasets over 20 training epochs. Additional results, including qualitative and task-specific examples as well as failure cases, are included in Appendix~\ref{app:more_qualitative},\ref{app:more_control_task},\ref{app:failure}.

\textbf{Convergence}. We assess how rapidly ControlNet, T2I-Adapter, LoRA, and COFT converge on the L2F task, utilizing both numerical and visual analyses. As our evaluation metric, we calculate the average $\ell_2$ norm between the locations of control and predicted face landmarks. Figure~\ref{fig:face_convergence} displays the progression of face landmark error as each model is finetuned across varying numbers of epochs. It is clear that COFT and OFT both exhibit notably accelerated convergence. For instance, LoRA requires 20 epochs to attain the performance level that OFT reaches at epoch 8. It is important to point out that both OFT and COFT have roughly the same parameter count as LoRA (and substantially fewer than ControlNet), yet they converge with far greater efficiency compared to prior methods. The rapid convergence behavior of OFT is further substantiated in Figure~\ref{fig:teaser}, where the rightmost example demonstrates OFT’s superior data efficiency relative to ControlNet and LoRA; specifically, OFT achieves satisfactory convergence with just 5\% of the ADE20K dataset. For qualitative analysis, we concentrate on OFT, COFT, and ControlNet since the latter produces landmark errors closest to those of our methods. As illustrated in Figure~\ref{fig:face_convergence_qual}, both OFT and COFT not only converge stably but also progressively synchronize the generated face pose with the guiding landmarks. In contrast, ControlNet's ability to follow control signals abruptly appears after the 8-th epoch, corresponding with the sudden convergence effect described in \cite{zhang2023adding}. This stable and consistent convergence pattern of OFT greatly facilitates its practical application, making the training trajectory smoother and more predictable and, as a result, simplifying the process of tuning OFT's hyperparameters.

\textbf{Quantitative comparison}. To assess the effectiveness of controllable generation, we propose a metric based on control consistency. This approach involves extracting the control signal from the generated output and directly comparing it to the initial input control signal. For the C2I task, we utilize IoU and F1 score, whereas for S2I we measure mean IoU as well as mean and overall accuracy. In the case of the L2F task, we calculate the average $\ell_2$ distance between the control landmarks and their predicted counterparts. Appendix~\ref{app:settings} provides further details regarding the evaluation metrics. Each method is evaluated using its optimal hyperparameter configuration. According to the results in Table~\ref{table:control_gen}, OFT and COFT demonstrate superior and more precise control than alternative approaches. Notably, adapter-based models such as T2I-Adapter and ControlNet exhibit slower convergence rates and ultimately achieve lower performance. LoRA, when compared with ControlNet, achieves higher performance in the S2I task but falls short in the C2I and L2F tasks. Overall, our findings suggest that LoRA’s maximum attainable performance remains limited, regardless of careful hyperparameter tuning. In contrast, empirical evidence indicates that the OFT framework continues to improve as the number of trainable parameters increases, suggesting that its performance has not reached a plateau. We highlight that this quantitative evaluation for controllable generation represents an important contribution of our work, enabling robust measurement of control fidelity in finetuned models, subject to the accuracy of the baseline segmentation or detection models employed.

\textbf{Qualitative comparison}. We further conduct a qualitative evaluation of OFT and COFT alongside leading methods such as ControlNet, T2I-Adapter, and LoRA. The samples displayed in Figure~\ref{fig:control_final}, generated at random, illustrate that both OFT and COFT not only produce highly realistic and high-quality images, but also demonstrate superior control capabilities. For the S2I task, LoRA is unable to generate outputs consistent with the provided segmentation map, whereas ControlNet, OFT, and COFT effectively respect the input constraints. Remarkably, OFT and COFT achieve visually richer details and more coherent geometric configurations compared to ControlNet, while utilizing significantly fewer model parameters. In the C2I setting, OFT and COFT both synthesize semantically consistent images from sparse Canny edge inputs, whereas T2I-Adapter and LoRA fall short in this regard. For the L2F task, our approach achieves the most precise pose-guided facial generation, even in the presence of challenging facial angles. Across all three control tasks, OFT and COFT consistently outperform state-of-the-art baselines in image quality, underlining the effectiveness of our OFT framework for controllable image synthesis. To broaden our qualitative assessment, additional visual results for all three tasks are available in Appendix~\ref{app:control_exp}. Furthermore, we showcase OFT's adaptability to a wider range of control problems—including dense pose-to-human body, sketch-to-image, and depth-to-image conversions—in Appendix~\ref{app:more_control_task}.

\section{Concluding Remarks and Open Problems}

Drawing inspiration from the role of angular relationships among neurons in shaping visual semantics, we introduce orthogonal finetuning, a straightforward and effective method for refining text-to-image diffusion models. Our approach is designed to address two particular applications: subject-driven generation and controllable generation. Relative to prior approaches, OFT achieves superior control and greater stability during finetuning while utilizing a smaller parameter footprint. Significantly, OFT maintains the same inference efficiency, making it suitable for practical deployment.

Nonetheless, OFT opens several avenues for further research. To enforce orthogonality, OFT employs Cayley parametrization, which necessitates computing a matrix inverse; this requirement imposes certain scalability constraints. While our adoption of block diagonal parametrization helps mitigate this issue, identifying differentiable techniques to accelerate the matrix inversion process remains unresolved. Furthermore, OFT presents intriguing opportunities for compositionality: since orthogonal matrices obtained from separate OFT finetuning sessions can be multiplied to yield another orthogonal matrix, investigating whether this process retains the specialized knowledge from each downstream task is a compelling topic for future work. Lastly, the parameter efficiency achieved by OFT hinges on the block diagonal formulation, which inherently imposes extra biases and curbs flexibility. Identifying strategies to enhance parameter efficiency while reducing such biases constitutes another central open question.